%\usepackage{graphicx}% já está incluido na classe
%\usepackage[utf8]{inputenc} % obsoleto, desnecessário

\usepackage[misc,geometry]{ifsym} 
\usepackage{pgfplots} 
\pgfplotsset{compat=1.17}

%%%%  %\usepackage{fontspec}

%% Os problemas com a classe eram originados por carregarem ambos os
%% pacotes fontenc e fontspec simultaneamente. Agora a classe detecta
%% qual engine está em uso e se for o pdflates então carrega o
%% fontenc. Se forem o xelatex ou lualatex então carrega o
%% fontspec. Eu particularemnte prefiro o lualatex ou o xetex


% \usepackage{fontawesome} %%% incluido na classe. Desnecessário
% chamá-lo aqui


%%%% \usepackage{academicons}
%%%% outro encrenqueiro que se tornou desnecessario. Deste typeface só
%%%% se usava o glifo do orcid e o codigo interno bombava.
%%%% Substitui pelo pacote orcidlink (carregado internamete na classe)
%%%% e consertei o codigo na classe.


%\usepackage{color} % carregada dentro da classe
%\usepackage{hyperref} % carregada dentro da classe

\usepackage{aas_macros}
\usepackage[bottom]{footmisc}

%\usepackage{supertabular}
%\usepackage{multicol}
%\usepackage{multirow}
%% O pacote tabularray é muito melhor para
% a criacao de tabela complexas e/ou loooongas; tudo de modo simples.
%% Torna o uso de multicol e multirow desnecessarios

\usepackage{tabularray}


\usepackage{afterpage}
\usepackage{url}
\usepackage{pifont}

\usepackage{booktabs}
\usepackage{float}

\setcitestyle{square}


\definecolor{engtitle}{rgb}{0.5,0.5,0.5}
\definecolor{orcidlogo}{rgb}{0.37,0.48,0.13}
\definecolor{unilogo}{rgb}{0.16, 0.26, 0.58}
\definecolor{maillogo}{rgb}{0.58, 0.16, 0.26}
\definecolor{darkblue}{rgb}{0.0,0.0,0.0}
\hypersetup{colorlinks,breaklinks,
            linkcolor=darkblue,urlcolor=darkblue,
            anchorcolor=darkblue,citecolor=darkblue}
%\hypersetup{colorlinks,citecolor=blue,linkcolor=blue,urlcolor=blue}

%%%%%%% IMPORTANT: We disable hyperlinks by default with this line, to avoid the error "\pdfendlink ended up in different nesting level" while writing.
%\hypersetup{draft}
